\documentclass[a4paper,11pt]{article}
% Подключаем все нужные пакеты
%%% Локализация и шрифты
\usepackage{asymptote}
\usepackage[utf8]{inputenc}          % Кодировка UTF-8
\usepackage[T2A]{fontenc}            % Русская кодировка
\usepackage[english,russian]{babel} % Локализация
\usepackage{cmap}                    % Поддержка поиска в PDF
\usepackage{tikz}
\usepackage{xstring}
\usepackage[left=20mm, top=10mm, right=20mm, bottom=10mm, nohead, nofoot]{geometry}
\renewcommand{\baselinestretch}{1.0}% Межстрочный интервал
\usepackage{ragged2e}               % Выравнивание текста
\usepackage{microtype}              % Улучшение качества вывода текста
\usepackage{multicol}
\usepackage{tcolorbox}
\tcbuselibrary{listingsutf8, breakable, skins}
\justifying
\sloppy
\tolerance=500
\hyphenpenalty=10000
\emergencystretch=3em

\usepackage[normalem]{ulem}
\pagestyle{empty} \textwidth=19.0cm \oddsidemargin=-1.3cm
\textheight=26cm \topmargin=-3.0cm
\tcbuselibrary{listingsutf8}
\usepackage{pgfplots}
\renewcommand{\rmdefault}{cmss}
\usepgfplotslibrary{patchplots}
\pgfplotsset{compat=1.18}
\usetikzlibrary{fillbetween}
\usetikzlibrary{patterns}
% Создаём новый счётчик задач

% определение
\newcounter{maindf}
\newcounter{subdf}[maindf]

\newcommand{\dfsub}{\arabic{maindf}.\arabic{subdf}}

% теорема
\newcounter{mainth}
\newcounter{subth}[mainth]

\newcommand{\thsub}{\arabic{mainth}.\arabic{subth}}

% утверждение
\newcounter{mainpr}
\newcounter{subpr}[mainpr]

\newcommand{\propsub}{\arabic{mainpr}.\arabic{subpr}}

% пример
\newcounter{mainex}
\newcounter{subex}[mainex]

\newcommand{\exsub}{\arabic{mainex}.\arabic{subex}}

% задача
\newcounter{mainta}
\newcounter{subta}[mainta]

\newcommand{\tasub}{\arabic{mainta}.\arabic{subta}}

\newcounter{mainle}
\newcounter{suble}[mainle]

\newcommand{\lesub}{\arabic{mainle}.\arabic{suble}}

\renewcommand{\familydefault}{\ttdefault}

% Определяем стили блоков
\usepackage{xcolor}
% myrose = #FFB5C1
% myred = #FF6464
% mybrown = #D2B48C
% mycyan = #87CEEB
% myorange = #FFB000
% mypurple = #D8BFD8
% mygreen = #14DE14
% --- Однострочное создание блоков \problem, \solution, \answer ---
\definecolor{myyellow}{HTML}{737700} % Жёлтый
\definecolor{mygreen}{HTML}{03810b}  % Голубой
\definecolor{mycyan}{HTML}{87CEEB}  % Голубой

\newtcolorbox{SectionBox}{
  breakable,
  enhanced,
  colback=yellow,
  colframe=myyellow,
  fonttitle=\ttfamily,
}

\newcommand{\sect}[1]{%
  \refstepcounter{subdf}%
  \begin{SectionBox}%
    \begin{center}
      #1
    \end{center}%
  \end{SectionBox}%
}

\newtcolorbox{SubsectionBox}{
  breakable,
  enhanced,
  colback=green,
  colframe=mygreen,
  fonttitle=\ttfamily,
}

\newcommand{\subsect}[1]{%
  \refstepcounter{subdf}%
  \begin{SubsectionBox}%
    \begin{center}
      #1
    \end{center}%
  \end{SubsectionBox}%
}

\newtcolorbox{SubsubsectionBox}{
  breakable,
  enhanced,
  colback=mycyan,
  colframe=cyan,
  fonttitle=\ttfamily,
}

\newcommand{\subsubsect}[1]{%
  \refstepcounter{subdf}%
  \begin{SubsubsectionBox}%
    \begin{center}
      #1
    \end{center}%
  \end{SubsubsectionBox}%
}

% Команда для векторов
\newcommand{\vect}[1]{\mathbf{#1}}
\usepackage[upint]{stix2}

%%% Колонтитулы
\usepackage{fancyhdr}
\pagestyle{fancy}
\renewcommand{\headrulewidth}{0.3mm} % Толщина линии в колонтитулах

\rhead{Динамическая оптимизация в экономике и финансах}
\lhead{ФЭН ВШЭ, 2026}
\headsep=25pt
\footskip=15mm
\textheight=700pt
\headheight=35pt % Возможно, понадобится увеличить до ~30pt. Проверьте log-файл.
\parindent=0mm
%%% Математика
\usepackage{amsmath, amsfonts, amssymb, amsthm, mathtools} % Базовые математические пакеты
\usepackage{icomma}   
\usepackage{euscript}               % Шрифт Евклид
\usepackage{mathrsfs}               % Красивый матшрифт
\usepackage{cancel}                 % Зачёркивание в формулах
\usepackage{diagbox}
\usepackage[unicode, colorlinks, urlcolor=darkblue, linkcolor=violet]{hyperref}
\usetikzlibrary{matrix}
% Замены часто употребляемых шрифтов
\newcommand{\mb}[1]{\mathbb{#1}}
\newcommand{\mc}[1]{\mathcal{#1}}
\newcommand{\ms}[1]{\mathscr{#1}}

\newcommand{\cov}{\operatorname{cov}}
\newcommand{\Var}{\operatorname{Var}}
\newcommand{\corr}{\operatorname{corr}}
\newcommand{\const}{\operatorname{const}}
\newcommand{\Be}{\operatorname{Be}}
\newcommand{\Bi}{\operatorname{Bi}}
\newcommand{\Poi}{\operatorname{Poi}}
\newcommand{\Exp}{\operatorname{Exp}}
\newcommand{\Geo}{\operatorname{Geo}}
\newcommand{\bias}{\operatorname{bias}}
\newcommand{\MSE}{\operatorname{MSE}}
\newcommand{\sign}{\operatorname{sign}}
\newcommand{\plim}{\operatorname{plim}}
\newcommand{\asi}{\overset{asy}{\sim}}
\newcommand{\argmax}[1]{\underset{#1}{\operatorname{argmax}}}
\begin{document}
\begin{center}
    \huge ДИНАМИЧЕСКАЯ ОПТИМИЗАЦИЯ В ЭКОНОМИКЕ И ФИНАНСАХ\\[10cm]
\end{center}
\begin{center}
    Конспекты лекций Пильника Н.Е.\\[2cm]

    Конспект подготовлен студентом ЭАДа Александром Благий\\[10cm]

    Факультет Экономических наук, 2026
\end{center}
\newpage
\tableofcontents
\newpage
\begin{center}
    \huge{Оптимизаиция в дискретном времени}\\[1cm]
    \large{Статическая оптимизация}\\[1cm]
\end{center}
Общая задача: $f(x) \to \max,\ x\in X \subseteq \mb{R}^n$.

Решение: $x^*\in X^*$

Тогда $x^*\ -$ точка максимума, а $f(x^*)\ -$ максимум функции $f$ на множестве $X$.\\[5mm]
Случаи отстутствия максимума:
\begin{enumerate}
    \item $X = \varnothing$;
    \item $X \neq \varnothing$, но $f(x)$ не ограничена сверху на $X$;
    \item $X \neq \varnothing$, $f(x)$ ограничена сверху, но точная верхняя грань не достигается.
\end{enumerate}
\addcontentsline{toc}{section}{Оптимизация в дискретном времени}
Если $g\ -$ монотонно возрастающая, то $g(f(x)) \to \max$ даёт то же $X^*$.

Любая точка экстремума является:
\begin{enumerate}
    \item максимумом/минимумом;
    \item глобальной/локальной;
    \item строгой/нестрогой;
    \item внутренней/граничной;
    \item условной/безусловной.
\end{enumerate}
Если $\exists \varepsilon \ \forall y\in \varepsilon(x)\ (y\neq x)$ и $f(y) > f(x)$, то $x\ -$ точка локального (строгого) максимума.

\begin{center}
    Безусловный экстремум
\end{center}
\[f(x) \to \max \quad f\ - \text{дважды дифференцируемая}\]
Необходимое условие: $\nabla f(x_0) = 0$, где $\nabla f\ -$ градиент функции $f$.

Достаточное условие: если $H \prec 0$, где
\[H = \begin{pmatrix}
\dfrac{\partial^2 f}{\partial x_1^2} & \cdots & \dfrac{\partial^2 f}{\partial x_1 \partial x_n} \\[0.3cm]
\vdots & \ddots & \vdots \\[0.3cm]
\dfrac{\partial^2 f}{\partial x_n \partial x_1} & \cdots & \dfrac{\partial^2 f}{\partial x_n^2}
\end{pmatrix},\] 
то $\max$.

(то есть матрица Гессе отрицательно определена).
\renewcommand{\thletter}{Критерий Сильвестра}
\theorem{
\begin{enumerate}
    \item $\Delta_1 > 0,\dots,  \Delta_n > 0 \implies \min$;
    \item $\Delta_1 < 0, \Delta_2 > 0, \Delta_3 < 0,\ldots \implies \max$;
    \item $\exists$ отрицательный чётный минор $\implies$ не экстремум;
    \item Где-то стоят нули $\implies$ требуется дополнительное исследование.
\end{enumerate}
}

\begin{center}
    Условный экстремум с ограничениями типа равенства
\end{center}
\[\begin{cases}
    f(x) \to \max & x\in\mb{R}^n\\
    g(x) = b & b\in\mb{R}^m
\end{cases}\]
где $f,\ g$ дважды дифференцируемы, $m < n$, а $X = \{x\mid g(x) = b\}$

Условие Якоби выполнено, если $\operatorname{rk}\, J_g(x) = m$, где $J_g(x)\ -$ матрица Якоби функции $g$.\\[5mm]

В точках, где это условие не выполнено, лагранжиан не работает.

Для $\lambda \in\mb{R}^m$ введём функцию $\mc{L}(x,\lambda) = f(x) + \lambda^T(g(x) - b)$

Необходимое условие:
\[\begin{cases}
    \dfrac{\partial \mc{L}}{\partial x_1} = \cdots = \dfrac{\partial\mc{L}}{\partial x_n} = 0;\\[3mm]
    \dfrac{\partial \mc{L}}{\partial \lambda_1} = \cdots = \dfrac{\partial\mc{L}}{\partial \lambda_m} = 0.
\end{cases}\]

Составим $\widehat{H} = \begin{pmatrix}
    0 & J_g(x)\\
    J_g(x)^T & H
\end{pmatrix}\ -$ матрицу Гессе лагранжиана $\mc{L}$.

Смотрим на миноры порядка от $2m+1$ до $n + m$.

Варианты:
\begin{enumerate}
    \item все миноры положительны $\implies$ локальный минимум;
    \item знаки миноров чередуются, начиная с минуса $\implies$ локальный максимум;
    \item $\exists$ отрицательные миноры на чётном месте $\implies$ не экстремум;
    \item где-то стоят нули $\implies$ требуется дополнительное исследование.
\end{enumerate}
\begin{center}
    Условный экстремум с ограничениями типа неравенства
\end{center}
\[\begin{cases}
    f(x) \to \max & x\in\mb{R}^n\\
    g(x) \leq b & b\in\mb{R}^m
\end{cases}\]
где $f,\ g$ дважды дифференцируемы, $m$ может быть больше $n$, а $X = \{x\mid g(x) \leq b\}$

Для $\lambda \in\mb{R}^m$ введём функцию $\mc{L}(x,\lambda) = f(x) + \lambda^T(b - g(x))$

Условия Каруша-Куна-Таккера (они же дополяющей жесткости):
\[\begin{array}{l}
    \forall j = 1,\dots, m\ \lambda_j \geq 0,\\
    b_j - g_j(x) \geq 0,\\
    \lambda_j(b_j - g_j(x)) = 0.
\end{array}\]

\addcontentsline{toc}{section}{Вариационное исчисление}




\addcontentsline{toc}{section}{Оптимальное управление}




\addcontentsline{toc}{section}{Динамическое программирование}
\end{document}